\section{Open Questions and Theoretical Boundaries}
\label{app:open-questions}

This appendix addresses theoretical questions raised during independent
review of the derivation.  For each question we provide the best answer
currently available, state honest limitations, and identify the
resolving test where one exists.

% -------------------------------------------------------------------------
\subsection{Three Layers of the Framework}
\label{app:oq:layers}

The paper can be read as deriving a reasoning system from nothingness, but
it is useful to separate three layers that play different logical roles.
We name them L0, L1, and L2.

\paragraph{L0 (Bottom State).}
L0 is absolute nothingness $\noth$, the unfalsifiable starting point.  In
L0 there are no distinctions, no carrier, and no tests, meaning
$\Del = \varnothing$ and $\Carrier = \varnothing$ (\step{0}).

\paragraph{L1 (Meaning Operator).}
L1 is Axiom~A0 (Witnessability) together with its forced consequences.
This layer contains the objects and constraints that follow from
witnessability alone, for example the carrier and finite slice
($\Carrier$, $\FinSlice$), self-delimitation ($\SdMap$), executability
(a total evaluator $\UTM$ and a cost functional), the truth quotient
$\TQ{\Ledger}$, entropic time $\Tim$, and gauge invariance $\gauge$.  Informally, L1 is the
operator that turns a would-be distinction into meaning by requiring a
finite witness.

\paragraph{L2 (Model-Class Commitments).}
L2 collects canonical or conservative commitments that are not forced by A0
alone, but define the paper's specific model class.  In this paper these
include commutativity of causally independent tests (\step{5}), prior-free
minimax optimization for test selection (\step{18}), and the closed-universe
assumption that forbids external delimiter or verifier oracles
(\step{2}, \step{19}, \step{20}).  Alternative choices are A0-compatible,
but they yield structurally analogous, different frameworks.

\paragraph{Coupling Law resolution.}
The scheduling commitment---which test to run next---was previously
classified as L2 (a modeling choice, since $\CanEnum$ tie-breaking was
one of several A0-compatible options).  The Separator-Indistinguishability
Coupling Law (\cref{thm:bellman}, Part~(b)) resolves this: the
question-relative efficiency $\EffQ{\tau}{W}$ is the canonical dimensionless
ratio constructible from the framework's own quantities, and $\PiCanon$
tie-breaking is forced by gauge invariance.  Together they make the
test-selection map $\OptTest{q}{W}$ a deterministic function of $(W, q)$,
upgrading the scheduling commitment from L2 (model class) to L1 (forced
by A0).

\begin{table}[ht]
\centering
\small
\caption{Three layers and their logical roles in \Opoch.}
\begin{tabular}{@{}lll@{}}
\toprule
\textbf{Layer} & \textbf{Name} & \textbf{Content} \\
\midrule
L0 & Bottom state & $\noth$, no admissible distinctions, no carrier, no tests \\
L1 & Meaning operator & A0 and its forced consequences (carrier, tests, truth, time, gauge) \\
L2 & Model class & Canonical commitments (commutativity, minimax, closed universe) \\
\bottomrule
\end{tabular}
\end{table}

% -------------------------------------------------------------------------
\subsection{Can the Self-Computing Machine Be Built?}
\label{app:oq:machine}

The system $\mathcal{S}$ is specified by \cref{sec:derivation}.  The question here is
constructibility in the strict sense, meaning whether an implementation can
be built from finitely specified, individually computable components.

\paragraph{Five-component constructibility argument.}
A full implementation decomposes into five components, each with a finite
witness procedure:
\begin{enumerate}[nosep]
  \item \textbf{Total evaluator on $\Del^*$.}  A universal evaluator $\UTM$
        that executes any endogenous test description $\tau$ on any input
        $x$ and returns a well-defined outcome, as fixed by the
        executability substrate (\step{3}).
  \item \textbf{Self-delimiting parser.}  A parser for the prefix-free
        interface induced by $\SdMap$ (\step{2}), which decodes
        $1^{|s|}0s$ in at most $2|s|+1$ bit-reads.
  \item \textbf{Bellman minimax solver on finite domains.}  The Bellman
        recursion over the surviving answer set and the feasible test set
        is computable because the system operates on $\FinSlice$ and, for
        fixed remaining budget, the feasible test set is finite
        (\step{18}).
  \item \textbf{Receipt chain.}  A deterministic receipt system that
        serializes inputs and outputs, hashes them, and chains hashes.
        This uses standard SHA-256
        collision resistance as a modeling commitment.
  \item \textbf{$\Pit$-canonicalization.}  A canonicalization operator
        $\PiCanon$ that reduces gauge-equivalent representations to a
        canonical representative before hashing and comparison.  This
        component is defined in Definition~\ref{def:pi-canon}.
\end{enumerate}
Since each component is individually computable and their composition is
finite, the self-computing machine is constructible in principle.

\paragraph{Practical limitation.}
The Bellman minimax recursion has worst-case complexity exponential in
$|\TQ{\Ledger}|$ (the number of truth classes).  For large problem
domains, heuristic approximations or domain-specific pruning are
needed.  The framework accommodates this: any pruned solver that
cannot reach $\UNIQUE$ within budget simply outputs $\OMEGA$ with
the frontier witness, which is the honest response.

\paragraph{Status.}
\textbf{Constructible in principle, with exponential worst-case cost.}
The three-point audit (\S\ref{sec:verification}) demonstrates internal
consistency on the derivation's 34 steps.

% -------------------------------------------------------------------------
\subsection{Quantum Mechanics and Non-Commuting Observables}
\label{app:oq:quantum}

Non-commutativity arises in two distinct ways.  The first is already
internal to the framework, while the second requires a genuine extension.

\paragraph{(a) Resource-induced noncommutativity (already covered).}
Even when tests are pure functions, the \emph{set of affordable tests}
depends on prior spending.  If running $\tau_1$ consumes budget, then the
subsequent feasible set $\Del(\Tim)$ available for $\tau_2$ can change.
This produces an order effect in policy space, but it is handled by the
Bellman loop, because the value function explicitly conditions on remaining
budget (\step{18}).

\paragraph{(b) Disturbance-induced noncommutativity (requires extension).}
In quantum measurement, the outcome statistics of $\tau_2$ can depend on
whether $\tau_1$ was performed, because performing $\tau_1$ changes the
state.  This is an order effect at the level of the physical test
instrument, not merely at the level of budget.

\paragraph{Scope of the restriction.}
Case (b) is excluded by the current model class, which assumes that tests
do not disturb the underlying state, and therefore that causally
independent tests commute at the quotient level (Diamond Law, \step{5}).
This assumption propagates to the multiset ledger (\step{5}) and to
unconditional commutativity in the eraser algebra (\step{7}).  The
remainder of the derivation, including carrier, syntax, executability,
time, gauge, separator geometry, Bellman dynamics, and self-hosting, does
not require commutativity and can be carried over.

\paragraph{Extension path.}
The extension is needed only for case (b).  Replace the multiset ledger
(\step{5}) with an ordered (or partially ordered) ledger, and refine tests
from pure maps to instruments $\Instrument$ that return both an outcome and
a state update.  In such a ledger, the truth quotient becomes
order-dependent, and the Diamond Law applies only to commuting pairs:
\[
  e_{\tau_1} \circ e_{\tau_2} = e_{\tau_2} \circ e_{\tau_1}
  \quad\text{iff}\quad [\tau_1, \tau_2] = 0.
\]

\paragraph{Status.}
\textbf{Case (a) is handled, case (b) is not.}  Resource-induced order
effects are already captured by budget conditioning in the Bellman
recursion.  Disturbance-induced noncommutativity requires the
ordered-ledger, instrumented-test extension sketched above.

% -------------------------------------------------------------------------
\subsection{Continuous Spaces and Infinite Carriers}
\label{app:oq:continuous}

The derivation operates on $\FinSlice \subset \Carrier$ (\step{1}).
Extending the framework to continuous state spaces requires care, because
naively importing differential entropy introduces coordinate dependence.

\paragraph{Correct invariant replacement.}
The finite framework measures uncertainty in terms of truth-quotient
classes.  Accordingly, in the continuous extension one should work with a
$\Pit$-fixed measure on truth classes, not with coordinate densities.
Concretely, for a surviving set $W$ of truth-classes, replace:
\begin{itemize}[nosep]
  \item $|W|$ with $\mu_{\Pit}(W)$, the counting measure on
        truth-quotient classes.
  \item $\log |W|$ with the $\Pit$-fixed entropy
        $H_{\Pit}(W) = -\sum_{w \in W} \mu_{\Pit}(w)\,\log \mu_{\Pit}(w)$,
        which is defined on truth classes and is gauge-invariant.
\end{itemize}
By contrast, differential entropy
$h(W) = -\int p\,\log p\,d\mu$ depends on the choice of coordinates,
and is therefore not gauge-invariant.

\paragraph{Relative-entropy formulation.}
The correct continuous extension uses relative entropy (KL divergence)
against the uniform measure on truth classes.  This preserves gauge
invariance and aligns with the system's use of class-level ambiguity.

\paragraph{Status.}
\textbf{The finite-carrier framework extends naturally.}  Rigorous
treatment requires measure-theoretic machinery (Borel $\sigma$-algebras,
measurable selection theorems) beyond the current paper's scope.  The
finite-slice derivation is the foundational case; the continuous
extension is deferred to future work.

% -------------------------------------------------------------------------
\subsection{G\"{o}del's Incompleteness and Self-Hosting}
\label{app:oq:godel}

G\"{o}del's second incompleteness theorem \citep{godel1931} states that no
consistent formal system of sufficient strength can prove its own
consistency.  The self-hosting property (\step{19}) can be misread as such
a proof, but it is not.

\paragraph{Resolution.}
Self-hosting is the claim that $\FixPt(\mathcal{S}) = \mathcal{S}$,
meaning the system can verify each derivation step by encoding it as a
compiled finite contract, serializing it, and running itself on that
contract to obtain $\UNIQUE$ (\cref{thm:self-hosting}).  In this sense,
the fixed point is a $\Pit$-serialize fixed point, because
$\Ser(\cdot)$ produces a finite code and $\PiCanon$ removes gauge
redundancy before receipt formation.

This is categorically distinct from a consistency proof.  The system checks
object-level claims of the form, ``\step{$i$} is the unique A0-compliant
continuation of \step{0} through \step{$i{-}1$},'' which are checkable.
It does not establish the meta-level claim, ``the system never produces
incorrect output,'' which would amount to a global consistency guarantee.
G\"{o}del blocks the latter, not the former.

\begin{enumerate}[nosep]
  \item \textbf{Object-level:} ``\step{$i$} is the unique A0-compliant
        continuation'' (checkable by replay on a finite contract).
  \item \textbf{Meta-level:} ``the system never produces incorrect output''
        (not checkable in the sense prohibited by G\"{o}del).
\end{enumerate}

\paragraph{Status.}
\textbf{No contradiction with G\"{o}del.}  Self-hosting is object-level
verification, not meta-consistency proof.  The system can verify its
derivation steps but cannot prove its own consistency.

% -------------------------------------------------------------------------
\subsection{The Halting Problem and the Total Evaluator}
\label{app:oq:halting}

The halting problem establishes that no algorithm can decide whether an
arbitrary Turing machine halts.  Step~3 introduces an evaluator $\UTM$,
and the framework's operational semantics require that admissible tests
terminate within the allocated budget.

\paragraph{Resolution.}
The evaluator operates on budget-total tests over the finite slice.
Write $\Del_{\leq B}$ for the set of tests that halt within budget $B$ on
every input in $\FinSlice$, equivalently the set of tests that are
admissible at budget $B$ (written $\Adm_{\leq B}(\tau)$).

\begin{enumerate}[nosep]
  \item \textbf{Budget-totality on $\FinSlice$.}  A test $\tau$ is
        \emph{budget-total} if $\UTM(\tau, x)$ halts within $B$ steps for
        every $x \in \FinSlice$.  Because $\FinSlice$ is finite (\step{1})
        and $B$ is fixed, this property is decidable by bounded simulation.
  \item \textbf{TIMEOUT as first-class outcome.}  Tests that do not halt
        within budget return $\mathrm{TIMEOUT} \in A$, which is a
        well-defined outcome in the test's outcome alphabet---not an
        undefined state or a semantic gap.  The evaluator is therefore
        \emph{total by construction}: every input produces a well-defined
        output within finite time, where $\mathrm{TIMEOUT}$ is one such
        output.
  \item \textbf{Constructive admissibility filtering.}  Enumerate candidate
        programs by $\CanEnum$ (\step{4}), execute each on all
        $x \in \FinSlice$ for at most $B$ steps, and admit those in
        $\Del_{\leq B}$.  Tests that exceed budget on some $x$ are excluded
        from the feasible set.
\end{enumerate}

The key insight is that A0's finiteness requirement (every witness must
terminate within finite budget) converts the undecidable question ``does
$\tau$ halt on all inputs?'' into the decidable question ``does $\tau$
halt on all inputs in $\FinSlice$ within budget $B$?''  This is not an
approximation or workaround---it is the correct formalization of A0's
operational semantics: only budget-feasible tests are admissible, and
budget-feasibility is decidable.

\paragraph{Status.}
\textbf{No conflict with the halting problem.}  The framework operates
with budget-totality on the finite slice, which is decidable.
$\mathrm{TIMEOUT}$ is a first-class outcome, not an error or gap.
Tests that might not halt on unbounded inputs are simply infeasible
(cost exceeds budget) and are excluded from $\Del(\Tim)$.

% -------------------------------------------------------------------------
\subsection{What Does $\OMEGA$ Look Like Concretely?}
\label{app:oq:omega}

We formalize the information returned in the $\OMEGA$ case.

\begin{definition}[$\OMEGA$ Payload]
\label{def:omega-payload}
The $\OMEGA$ output is a 5-tuple $(\mathcal{A}_S, \tau_\OMEGA, c(\tau_\OMEGA),
B_{\mathrm{rem}}, w_{\min})$ where:
\begin{itemize}[nosep]
  \item $\mathcal{A}_S \subset \mathcal{A}$ is the surviving answer set,
        with $|\mathcal{A}_S| > 1$.
  \item $\tau_\OMEGA$ is the minimal separator test for $\mathcal{A}_S$.
  \item $c(\tau_\OMEGA)$ is the separator cost.
  \item $B_{\mathrm{rem}}$ is the remaining budget.
  \item $w_{\min}$ is the minimality witness, the Bellman search subtree
        certifying that no cheaper separator exists.
\end{itemize}
\end{definition}

\paragraph{Example.}
Query: ``Find collision-free paths for 10 robots on a $25 \times 25$
grid with 30 delivery jobs.''  After exhausting the test budget, the
system returns:
\[
  \OMEGA:\;\bigl(\mathcal{A}_S,\; \tau_\OMEGA,\; c(\tau_\OMEGA),\;
  B_{\mathrm{rem}},\; w_{\min}\bigr),
\]
where $|\mathcal{A}_S| = 3$, $\tau_\OMEGA$ is a corridor-conflict test, and
$c(\tau_\OMEGA) = 47$ while $B_{\mathrm{rem}} = 12$, so the separator is
not affordable.
\emph{Interpretation:} the system cannot distinguish the 3 candidates
with remaining budget, but tells the user \emph{exactly} what test would
resolve the ambiguity and how much additional budget is needed.  This is
qualitatively different from ``I don't know''---it is ``I don't know,
\emph{and here is how to find out.}''

% -------------------------------------------------------------------------
\subsection{Declared vs.\ Derived Assumptions---Honest Inventory}
\label{app:oq:assumptions}

\Cref{tab:assumptions} provides a complete inventory of assumptions
relevant to the system as formalized in this paper, with an explicit layer
classification matching \cref{app:oq:layers}.

\begin{table}[ht]
\centering
\footnotesize
\caption{Complete assumption inventory for the 34-step derivation.}
\label{tab:assumptions}
\begin{tabular}{@{}clcllp{3.8cm}@{}}
\toprule
\textbf{\#} & \textbf{Assumption} & \textbf{Step} & \textbf{Layer} & \textbf{Status} & \textbf{Justification} \\
\midrule
1 & Witnessability (A0) & 0 & L1 & Axiom & Foundational, the only law \\
2 & Finite slice $\FinSlice$ & 1 & L1 & Derived & Forced by A0 finiteness \\
3 & Totality of $\UTM$ on $\Del^*$ & 3 & L1 & Derived & Forced by executability layer \\
4 & Positive integer costs & 3 & L1 & Derived & Cost equals execution steps $\geq 1$ \\
5 & Determinism of system & 18 & L1 & Derived & Canonical enumeration and Bellman argmin \\
6 & Commutativity of tests & 5 & L2 & Modeling commitment & Pure, non-disturbing test model \\
7 & Prior-free optimization & 18 & L2 & Modeling commitment & Minimax as conservative prior-free policy \\
8 & Closed-universe assumption & 2, 19, 20 & L2 & Modeling commitment & No external delimiter or verifier \\
9 & $\Pit$-canonicalization uniqueness & --- & L1 & Derived & Definition~\ref{def:pi-canon} fixes a canonical representative \\
10 & Receipt hash collision resistance & --- & L2 & Modeling commitment & SHA-256 assumed collision-resistant \\
11 & Smoothness of efficiency landscape & 31 & L1 & Assumption & Required for complex structure $J$ \\
12 & Conductance $w = 1/d^2$ & 28 & L2 & Modeling commitment & Dimensional analysis; alternatives possible \\
\bottomrule
\end{tabular}
\end{table}

\noindent
The layer column makes the logical separation explicit.  L1 items are
meaning-operator consequences, and L2 items are model-class commitments.
Alternative L2 choices, for example an ordered ledger for disturbance or a
Bayesian objective for test selection, remain A0-compatible but change the
resulting system class.  Phase~H adds one further L2 commitment: the
conductance weights $w = 1/d^2$ in the Dirichlet form (\step{28}), which
follow from dimensional analysis but could be replaced by any
$w = f(d)$ with $f$ strictly decreasing.

% -------------------------------------------------------------------------
\subsection{Verification Challenges: Framework-Level Interpretations}
\label{app:oq:challenges}

The framework's vocabulary invites connections to open problems in physics
and mathematics.  \Cref{tab:challenges} organizes these connections into
three categories.

\begin{table}[ht]
\centering
\small
\caption{Open verification challenges and framework-level interpretations.
These are interpretive connections, not formal proofs within the
derivation.}
\label{tab:challenges}
\begin{tabular}{@{}lllp{5.1cm}@{}}
\toprule
\textbf{Challenge} & \textbf{Domain} & \textbf{Category} & \textbf{Framework Status} \\
\midrule
Arrow of Time & Thermodynamics & Structural & $\Del\Tim \geq 0$ forced (\step{8}), internal arrow of time \\
Renormalization & QFT & Contract instantiation & Gauge coequalization gives vocabulary (\step{12}), interpretive only \\
Probability Origins & Foundations & Contract instantiation & Test counting forces frequency ratios, not full measure theory \\
Measurement Problem & QM & Contract instantiation & Fiber shrinkage analogy (\step{8}), not a resolution \\
Hard Problem & Consciousness & Open boundary & Requires specifying new $\Del$ instruments and verifiers \\
P vs.\ NP & Complexity & Open boundary & Separator geometry (\step{16}) suggests a lens, no reduction exists \\
\bottomrule
\end{tabular}
\end{table}

\paragraph{Disclaimer.}
These entries are \emph{framework-level interpretations}---they map
external problems into the derivation's vocabulary.  None constitutes a
proof or resolution of the corresponding open problem.  They are
included to clarify the framework's scope and to identify where further
formalization could be productive.
